% lautools API manual: scan metadata, conversion of scans and the lautools command line tools.
% lautools-browser and the laupy tools are described in laupy_lautools_manual (maintained in laupy).
% Compile with:  tectonic lautools_api_manual.tex
\documentclass[11pt,a4paper]{article}

\usepackage{fontspec}
\usepackage[margin=2.3cm]{geometry}
\usepackage{xcolor}
\usepackage[most]{tcolorbox}
\usepackage{listings}
\usepackage{booktabs,tabularx,longtable,array}
\usepackage{enumitem}
\usepackage{titlesec}
\usepackage{fancyhdr}
\usepackage[hidelinks]{hyperref}

\definecolor{accent}{HTML}{1F5A96}
\definecolor{termbg}{HTML}{F4F6F8}
\definecolor{termframe}{HTML}{B8C2CC}
\definecolor{notebg}{HTML}{EAF2FB}
\definecolor{warnbg}{HTML}{FDF1E4}
\definecolor{warnframe}{HTML}{D9822B}

\hypersetup{colorlinks=true,linkcolor=accent,urlcolor=accent}
\setlist{itemsep=2pt,topsep=4pt}
\setlength{\parindent}{0pt}
\setlength{\parskip}{5pt}
\setlength{\emergencystretch}{3em}

\titleformat{\section}{\Large\bfseries\sffamily\color{accent}}{\thesection}{0.8em}{}
\titleformat{\subsection}{\large\bfseries\sffamily}{\thesubsection}{0.8em}{}
\titleformat{\subsubsection}{\normalsize\bfseries\sffamily}{\thesubsubsection}{0.8em}{}

\pagestyle{fancy}
\fancyhf{}
\fancyhead[L]{\small\sffamily lautools --- API and command line tools}
\fancyhead[R]{\small\sffamily\nouppercase{\leftmark}}
\fancyfoot[C]{\small\thepage}
\renewcommand{\headrulewidth}{0.4pt}

\lstset{
  basicstyle=\ttfamily\small,
  breaklines=true,
  breakatwhitespace=false,
  postbreak=\mbox{\textcolor{gray}{$\hookrightarrow$}\space},
  columns=fullflexible,
  keepspaces=true,
  showstringspaces=false,
  tabsize=4,
  upquote=true,
  language={},
}

\tcbset{terminalstyle/.style={listing only, colback=termbg, colframe=termframe, boxrule=0.4pt,
  arc=2pt, left=4pt, right=4pt, top=2pt, bottom=2pt, breakable}}
\newtcblisting{terminal}{terminalstyle}
\newtcblisting{terminalsmall}{terminalstyle, listing options={basicstyle=\ttfamily\footnotesize}}
\newtcolorbox{note}{colback=notebg, colframe=accent, boxrule=0.4pt, arc=2pt, left=6pt, right=6pt, breakable}
\newtcolorbox{warning}{colback=warnbg, colframe=warnframe, boxrule=0.4pt, arc=2pt, left=6pt, right=6pt, breakable}

\newcommand{\code}[1]{\texttt{\detokenize{#1}}}

\begin{document}
\sloppy

\begin{titlepage}
\centering
\vspace*{4cm}
{\Huge\bfseries\sffamily\color{accent} lautools\par}
\vspace{0.6cm}
{\Large\sffamily Python API and command line tools\par}
\vspace{0.4cm}
{\large\sffamily scan metadata, conversion of scans, processing tools\par}
\vspace{1.2cm}
{\large Manual\par}
\vfill
{\small Vojtěch Kulvait, Helmholtz-Zentrum Hereon\par}
{\small November 2026\par}
\end{titlepage}

\tableofcontents
\newpage

% ====================================================================
\section{Overview}

lautools is the tomography package of the laupy ecosystem. It has three parts:
\begin{itemize}
\item \textbf{lautools-browser}, the graphical front end of projects, recipes and pipelines; it is
  described with the laupy command line tools in the \emph{laupy \& lautools-browser user manual}
  (\code{laupy_lautools_manual.pdf}).
\item \textbf{the Python API} of scans: metadata of PETRA III scans, the conversion of scans to zarr
  and DEN, processing helpers --- this manual.
\item \textbf{command line tools} built on the API, used by the pipelines of the recipes --- this
  manual, \S\ref{sec:tools}.
\end{itemize}

\begin{longtable}{@{}p{3.1cm}p{11.7cm}@{}}
\toprule
\textbf{Module} & \textbf{Content}\\
\midrule
\endhead
\code{PETRA} & Scan metadata of PETRA III (P05, P07): HDF5 files and zarr stores parsed into one canonical \code{ScanDescription}; analyses (frames, current, experiment info); writers to zarr and HDF5 (\S\ref{sec:petra}).\\
\code{TOMOGRAPHY} & Functions on scan datasets independent of the instrument: the TIFF stack specification of a scan (\S\ref{sec:tomography}).\\
\code{NANOCT} & Scan datasets of P05 nanotomography from \code{LogScan.log} and the TIFF files (\S\ref{sec:nanoct}).\\
\code{COR} & Rotation centre estimation with a JSON cache.\\
\code{preprocess} & Frame preprocessing: hot pixels (median or zinger filter), large components.\\
\code{scripts} & The command line tools (\S\ref{sec:tools}).\\
\bottomrule
\end{longtable}

lautools uses denpy for the data formats: DEN files, zarr codecs, the parallel zip store
\code{ZARZS} and the TIFF stack conversion \code{denpy.TIFF}; see the \emph{denpy manual}. The division
of labour: denpy knows formats and has no measurement specific code; lautools knows scans and
instruments and builds the conversion specifications that denpy executes.

% ====================================================================
\section{Scan metadata: PETRA}\label{sec:petra}

\subsection{One canonical scan description}

The metadata of a PETRA III scan exist in several forms: NeXus HDF5 files written by the beamline
software in different years and beamlines, and the \code{params} group of zarr stores written by the
conversion tools. The layouts differ in details that carry no meaning:

\begin{longtable}{@{}p{3.2cm}p{11.6cm}@{}}
\toprule
\textbf{Layout} & \textbf{Typical form}\\
\midrule
\endhead
P05 2021 & Attributes as one-element object arrays, single strings as (1,) datasets, datasets in chunks of 1024.\\
P05 2025 & \code{NX_class} as plain string, scalar string datasets, empty (0,) logs, chunks of one element, unlimited first dimension.\\
P07 & (1,) strings, each dataset one chunk of its full length.\\
zarr \code{params} & Scalar datasets as group attributes, N-D datasets as arrays, plus the scan table.\\
\bottomrule
\end{longtable}

\code{lautools.PETRA} parses every form into one canonical \code{ScanDescription} and implements every
analysis once, on it. The canonical form follows the zarr data model: a tree of \code{ScanGroup}s,
each with
\begin{itemize}
\item \textbf{attributes}: JSON values: the HDF5 attributes and the scalar datasets; numpy values as
  Python values, bytes as \code{str}, one-element sequences as their element; an attribute wins
  over a scalar dataset of the same name;
\item \textbf{arrays}: numpy arrays with their dtype; strings as object arrays of \code{str};
\item \textbf{groups}.
\end{itemize}
A string array of one element is a string attribute (an artefact of the writing software, e.g.\
\code{entry/beamline/name} of 2021 files), except in logs --- \code{NXlog} groups or groups with
\code{time} and \code{value} --- whose arrays stay arrays whatever their length. A (1,) number stays an
array; \code{value()} gives its element.

\subsection{Reading: ScanDescriptionFactory}

\begin{terminal}
from lautools import PETRA
d = PETRA.ScanDescriptionFactory.fromH5("scan.h5")            # file name, h5py File or Group
d = PETRA.ScanDescriptionFactory.fromZarr("scan.zar")         # store name or zarr.Group; group params
d = PETRA.ScanDescriptionFactory.create(source)               # any of these, or a ScanDescription
\end{terminal}

\code{fromZarr} leaves out the scan table (attributes \code{columns}, \code{index}, \code{data} of
\code{params}) and takes the path of the HDF5 file of the scan from the root attributes (\code{h5}).

\begin{note}
\textbf{Lazy reading.} P05 files have about 150 groups and store their arrays in chunks of one
element, which are slow to read. From an HDF5 file, groups and arrays are read when first used, the
members of a group at a time; \code{getExperimentInfo} reads a handful of groups. The file has to stay
in place until the description is used; \code{d.load()} reads everything.
\end{note}

\subsection{ScanGroup}

\begin{longtable}{@{}p{4.4cm}p{10.4cm}@{}}
\toprule
\textbf{Member} & \textbf{Meaning}\\
\midrule
\endhead
\code{attrs}, \code{arrays}, \code{groups} & The members (dicts).\\
\code{value(path)} & The scalar at \code{path}, wherever it is stored: an attribute, or the element of a one-element array; \code{None} for an empty array; a longer array as it is.\\
\code{data(path)} & The values as a 1-D numpy array (an array raveled, an attribute as one element).\\
\code{array(path)}, \code{group(path)} & The array or group, or \code{None}.\\
\code{path in g}, \code{g[path]} & Groups, arrays and attributes.\\
\code{scalars()} & The values of the arrays and attributes of the group except \code{NX_class}.\\
\code{differences(other)} & List of differences (attributes, arrays with dtype and values, groups); \code{==}.\\
\bottomrule
\end{longtable}

Paths are relative to \code{/entry} (HDF5) or \code{params} (zarr): \code{d.root.value("beamline/name")}.

\subsection{ScanDescription}

\begin{longtable}{@{}>{\raggedright\arraybackslash}p{6.2cm}p{8.6cm}@{}}
\toprule
\textbf{Method} & \textbf{Result}\\
\midrule
\endhead
\code{scanDataset(includeCurrent,} \code{timeOffsetSec)} & DataFrame of the frames in acquisition order: \code{time}, \code{image_key} (0 projection, 1 flat, 2 dark), \code{image_file} and the logs of \code{scan/data} (\code{s_rot}, \code{s_stage_x}, \dots); with \code{includeCurrent} the beam current interpolated at the frame times. Index: UNIX time in ms.\\
\code{imageDataset(image_key,} \code{includeCurrent,} \code{includePixelShift)} & The frames of one key with \code{frame_ind}, optionally \code{pixel_shift} from \code{s_stage_x}.\\
\code{getExperimentInfo(} \code{overrideMagnification)} & Experiment, beamline, times, frame rates, setup (distance, energy), camera (pixel size, field of view), Fresnel number.\\
\code{beamCurrentDataset(} \code{timeOffsetSec)} & Beam current readings; \code{BeamCurrentError} if there are none.\\
\code{findInsertionEvents(...)} & Top-up injections in the current.\\
\code{getBasicDatasetInfo()} & The frame table as logged (ms times).\\
\code{rootAttributes(export=None)} & JSON root attributes of a zarr store of the scan: experiment info and \code{export}.\\
\code{writeZarrGroup(group,} \code{name="params",} \code{scanTable=True)} & Write the description (and the scan table) as a zarr group.\\
\code{writeH5(path, chunk=1)} & Write \code{/entry} of an HDF5 file in the 2025 P05 layout.\\
\code{load()}, \code{differences(other)}, \code{==} & Read everything; compare.\\
\bottomrule
\end{longtable}

\subsection{Module functions and migration from denpy}

The module functions of \code{denpy.PETRA} exist with the same signatures and results:
\code{scanDataset}, \code{imageDataset}, \code{getExperimentInfo}, \code{beamCurrentDataset},
\code{findInsertionEvents}, \code{getBasicDatasetInfo}, \code{BeamCurrentError},
\code{scanDatasetNanotomography}. They take an HDF5 file (name or h5py object), a zarr store with a
\code{params} group, or a \code{ScanDescription}. Migration is the import:

\begin{terminal}
from lautools import PETRA        # was: from denpy import PETRA
df = PETRA.scanDataset("scan.h5", includeCurrent=True)
df = PETRA.scanDataset("raw.zar", includeCurrent=True)       # the same from the converted store
\end{terminal}

\code{denpy.PETRA} stays as the legacy implementation; the lautools tests compare both on synthetic
scans and on real beamline files.

\subsection{Writing: zarr and HDF5}

\begin{terminal}
d = PETRA.ScanDescriptionFactory.fromH5("scan.h5")
root = zarr.open_group(store, mode="a")
d.writeZarrGroup(root)                        # group params; scanTable=df for a given table
d.writeH5("reconstructed.h5")                 # /entry in the 2025 layout
\end{terminal}

\code{writeH5} writes \code{NX_class} as plain string attributes, the other attributes as scalar
datasets (strings as variable-length ASCII, UTF-8 if needed), arrays with unlimited first dimension in
chunks of one element and no compression, as the beamline does. A 2025 file is reproduced exactly; an
older file becomes the same description in the 2025 layout. The syscripts tool
\code{zarrToBeamtimeH5.py} is \code{fromZarr(...).writeH5(...)}.

\subsection{Ground truth and tests}

\code{lautools/test_PETRA.py} checks: the border cases of the files (missing or empty current, empty
or scalar camera values, incomplete logs); that both P05 layouts give the same description; zarr round
trips (directory and zip store) and that every analysis gives the same result from zarr as from HDF5;
that \code{writeH5} reproduces the 2025 layout; and that every analysis gives the results of
\code{denpy.PETRA}. With real files:

\begin{terminal}
LAUTOOLS_PETRA_TEST_FILES=/path/to/h5_examples python -m pytest lautools/test_PETRA.py
\end{terminal}

% ====================================================================
\section{Scan datasets to conversion specifications: TOMOGRAPHY}\label{sec:tomography}

A \emph{scan dataset} is the DataFrame of the frames of a scan, as \code{PETRA.scanDataset} and
\code{NANOCT.scanDataset} create it: columns \code{image_key} and \code{image_file} (and
\code{file_exists} for nanoCT), rows in acquisition order.
\code{TOMOGRAPHY.scanTiffStackConversionSpecs} turns it into a TIFF stack specification of denpy
(see the denpy manual): one tiff stack and one array per image key.

\begin{terminal}
spec = TOMOGRAPHY.scanTiffStackConversionSpecs(df, output, rawDir=None, names=None,
          fillValue=None, frameShape=None, dtype=None, defaults=None, attributes=None)
\end{terminal}

\begin{longtable}{@{}p{2.6cm}p{12.2cm}@{}}
\toprule
\textbf{Argument} & \textbf{Meaning}\\
\midrule
\endhead
\code{output} & The output of the specification, e.g.\ \code{{"format": "zarr", "store": "zip", "path": "raw.zar"}}.\\
\code{rawDir} & \code{image_file} is relative to it (a leading \code{/} is ignored, as in P05 files); \code{None}: paths as they are (nanoCT: absolute).\\
\code{names} & Image key $\to$ name; default \code{SCAN_ARRAY_NAMES} \code{{2: "dar", 1: "ref", 0: "img"}}, the names the pipelines read. A key without rows gives a stack of 0 frames.\\
\code{fillValue} & \code{None}: a row without a file (empty name, \code{file_exists} false, not on disk) is an error; a number: such rows become \code{null} frames written as this value and listed in the report.\\
\code{frameShape}, \code{dtype} & Of the frames; default from the first existing file.\\
\code{defaults}, \code{attributes} & Array defaults (e.g.\ the codec) and root attributes.\\
\bottomrule
\end{longtable}

The result is an ordinary specification: extend it (\code{denpy.TIFF.extend_spec}), save it, run it
with \code{denpy.TIFF.execute} or \code{convertTiffStack}.

\begin{terminal}
from lautools import PETRA, TOMOGRAPHY
from denpy import TIFF
scanDescription = PETRA.ScanDescriptionFactory.fromH5("scan.h5")
df = scanDescription.scanDataset()
spec = TOMOGRAPHY.scanTiffStackConversionSpecs(df, {"format": "den", "path": "origin"}, rawDir="raw")
spec = TIFF.extend_spec(spec, "dark_correction.json")
report = TIFF.execute(TIFF.validate(spec))
\end{terminal}

% ====================================================================
\section{Nanotomography: NANOCT}\label{sec:nanoct}

\code{NANOCT.scanDataset(LogScan, imgDir=None, one_based=True, drop_missing=True,
check_block_order=True)} builds the scan dataset of a P05 nanoCT measurement from \code{LogScan.log}
and the raw TIFF files (\code{img_*}, \code{ref_*}, \code{dar_*}): log rows and files are matched by
blocks and counters; columns \code{time}, \code{image_key}, \code{image_file}, \code{file_exists},
\code{s_rot}, \code{current}, \code{log_counter}, \code{log_block}.
\code{scanDatasetWithForeignDarkFields} appends the flat fields of another measurement as dark
fields.

% ====================================================================
\section{Command line tools}\label{sec:tools}

\begin{longtable}{@{}>{\raggedright\arraybackslash\footnotesize}p{7.4cm}>{\raggedright\arraybackslash}p{7.4cm}@{}}
\toprule
\textbf{Command} & \textbf{Purpose}\\
\midrule
\endhead
\code{tiffScanDataToZarrForMicrotomography} & Zarr store or DEN files of a PETRA III microtomography scan; replaces \code{rawTiffStructureToZarr.py} and \code{rawTiffStructureToDen.py} (\S\ref{sec:micro}).\\
\code{tiffScanDataToZarrForNanotomography} & Zarr store of a nanoCT scan (\code{LogScan.log}).\\
\code{tiffScanInfoForMicrotomography}, \code{tiffScanInfoForNanotomography} & Per-frame statistics of the raw frames with the metadata.\\
\code{createWorkingDirectoryForMicrotomography}, \code{createWorkingDirectoryForNanotomography} & Working directories with the selected samples.\\
\code{createTickFile} & DEN file of angle, shift and current of the frames.\\
\code{currentCorrect} & Beam current correction after dark field correction.\\
\code{selectAngleSequenceForFBP} & Contiguous angular range for 180\textdegree{} filtered back projection.\\
\code{removeHotPixels} & Hot pixels and zingers of a DEN or zarr stack.\\
\code{binDenFile} & Binning of DEN files.\\
\code{lautools-browser} & The graphical front end (other manual).\\
\bottomrule
\end{longtable}

The tools that read scan metadata use \code{lautools.PETRA}, so they accept the HDF5 file of a scan.
Every tool prints its options with \code{--help}.

\subsection{tiffScanDataToZarrForMicrotomography}\label{sec:micro}

The conversion of a microtomography scan, and the worked example of the API:

\begin{enumerate}
\item metadata $\to$ scan dataset:\\ \code{scanDescription = PETRA.ScanDescriptionFactory.fromH5(h5)},\\
  \code{df = scanDescription.scanDataset()};
\item scan dataset $\to$ specification: \code{TOMOGRAPHY.scanTiffStackConversionSpecs(df, ...)},
  extended by \code{--spec-extension};
\item specification $\to$ arrays: \code{TIFF.execute(TIFF.validate(spec))};
\item metadata $\to$ the same store: \code{scanDescription.writeZarrGroup(group, scanTable=df)} in
  mode \code{"a"} (the zip is appended safely by ZARZS); DEN: \code{params.json}.
\end{enumerate}

\begin{terminalsmall}
tiffScanDataToZarrForMicrotomography INPUT.h5 OUTPUT [--raw-dir DIR] [--zip] [--v2]
    [--compression NAME] [--clevel L] [--den] [--dtype DTYPE] [--force] [-j THREADS]
    [--fill-missing VALUE] [--fix-corrupted-h5] [--spec-extension EXT.json ...]
    [--write-spec spec.json] [--report report.json] [--check-headers] [--dry-run] [--verbose]
\end{terminalsmall}

\begin{longtable}{@{}p{4cm}p{10.8cm}@{}}
\toprule
\textbf{Option} & \textbf{Meaning}\\
\midrule
\endhead
\code{OUTPUT} & Zarr store (directory, or zip with \code{--zip}; \code{.zar} recommended) or, with \code{--den}, the directory of \code{dar.den}, \code{ref.den}, \code{img.den}.\\
\code{--raw-dir} & Directory of the TIFF files; default the directory of the HDF5 file.\\
\code{--compression}, \code{--clevel} & Zarr codec (default blosc-zstd 5; htj2k 0 is lossless).\\
\code{--den} & DEN files instead of zarr; the directory may exist (e.g.\ the working directory) and its other files are not touched.\\
\code{--dtype} & dtype of dar, ref, img; default the TIFF dtype (uint16), \code{float32} as \code{rawTiffStructureToDen.py}.\\
\code{--force} & Replace existing output (for DEN: the files of the arrays).\\
\code{--fill-missing} & Frames listed in the HDF5 file without a TIFF file are written as this value (null frames, in the report); default an error.\\
\code{--fix-corrupted-h5} & If the HDF5 file lists no projections, \code{img} are the TIFF files of the raw directory that are not darks or flats.\\
\code{--spec-extension} & Virtual stacks and arrays added to the generated specification (repeatable).\\
\code{--write-spec}, \code{--report} & Save the specification and the report (JSON).\\
\code{--dry-run} & Validate and print what would be written.\\
\bottomrule
\end{longtable}

\subsubsection{Replacing the syscripts tools}

\begin{terminalsmall}
rawTiffStructureToZarr.py h5 h2k_ll.zar --zip --compression htj2k --clevel 0 --raw-dir raw --force
tiffScanDataToZarrForMicrotomography h5 h2k_ll.zar --zip --compression htj2k --clevel 0 --raw-dir raw --force

rawTiffStructureToDen.py h5 . --raw-dir raw --force
tiffScanDataToZarrForMicrotomography h5 . --den --dtype float32 --raw-dir raw --force
\end{terminalsmall}

The zarr stores are the same (tested): arrays, root attributes, and the \code{params} group as the same
scan description, with two improvements: single strings of 2021 files are strings, not one-element
arrays, and image codecs get no additional zstd compressor. The DEN files are bit identical to
\code{rawTiffStructureToDen.py}. Without \code{--dtype float32} the DEN files keep uint16, half the
size.

\subsubsection{Processing during the conversion}

A specification extension adds virtual stacks and arrays; an array of the same name replaces the
generated one, \code{null} removes it. Dark field correction with the median of the darks, writing the
median too:

\begin{terminalsmall}
{
  "virtualstacks": {"dar_median": {"stack": "dar", "reduce": "median", "dtype": "uint16"}},
  "arrays": {
    "dar_median": {"stack": "dar_median"},
    "img": {"stack": "img", "dtype": "uint16", "clip": true, "args": ["dar_median"],
            "lambda": "lambda x, d: x.astype('int32') - d"},
    "ref": {"stack": "ref", "dtype": "uint16", "clip": true, "args": ["dar_median"],
            "lambda": "lambda x, d: x.astype('int32') - d"}
  }
}
\end{terminalsmall}

Flat field corrected projections in float32 with the mean flat field, without writing the raw
projections:

\begin{terminalsmall}
{
  "virtualstacks": {"dar_mean": {"stack": "dar", "reduce": "mean"},
                    "ref_mean": {"stack": "ref", "reduce": "mean"}},
  "arrays": {
    "img": null,
    "img_ffc": {"stack": "img", "dtype": "float32", "args": ["dar_mean", "ref_mean"],
                "lambda": "lambda x, d, f: (x - d) / np.maximum(f - d, 1)"}
  }
}
\end{terminalsmall}

2$\times$2 binning of 5120$\times$3840 projections: \code{"lambda x: x.reshape(1920, 2, 2560, 2).mean(axis=(1, 3))"}.
Longer processing goes into a function of an importable module (on \code{PYTHONPATH}), called as
\code{f(frame, *args)} and returning the frame; e.g.\ a hot pixel filter around
\code{lautools.preprocess.remove_hot_pixels}, which returns the frame, the mask and a count:

\begin{terminalsmall}
# myfilters.py
from lautools import preprocess
def removeHotPixels(frame, iterations, filterSize):
    return preprocess.remove_hot_pixels(frame.astype("float32"), iterations, filterSize)[0]

"img": {"stack": "img", "dtype": "float32", "function": "myfilters:removeHotPixels", "args": [1, 5]}
\end{terminalsmall}

\begin{terminal}
tiffScanDataToZarrForMicrotomography h5 raw.zar --zip --compression htj2k --clevel 0 \
    --raw-dir raw --spec-extension dark_correction.json --dry-run
\end{terminal}

\subsubsection{The specification as a record}

\code{--write-spec} saves the specification that was run, with every frame listed; it can be run
again or modified with denpy's \code{convertTiffStack}, independently of the HDF5 file. A suggestion
for recipes (not enforced by any tool): keep the files of the original conversion of a scan together in
\code{WD/origin/} --- the specification, the report and the zarr store or DEN files --- so that they
can be kept or removed easily when the working directory gets populated:

\begin{terminal}
mkdir -p origin
tiffScanDataToZarrForMicrotomography h5 origin/raw.zar --zip --compression htj2k --clevel 0 \
    --raw-dir raw --write-spec origin/spec.json --report origin/report.json
\end{terminal}

\subsubsection{Output layout (zarr)}

\begin{longtable}{@{}p{3.8cm}p{11cm}@{}}
\toprule
\textbf{Entry} & \textbf{Content}\\
\midrule
\endhead
root attributes & Experiment info (\code{getExperimentInfo}) and \code{export} (dimensions, counts, dtype, codec, paths).\\
\code{dar}, \code{ref}, \code{img} & The frames, one per chunk, in acquisition order.\\
\code{params} & The scan description (the HDF5 entry) and the scan table (\code{columns}, \code{index}, \code{data}: pandas \code{orient="split"}).\\
\bottomrule
\end{longtable}

The metadata of a converted store are read back with \code{PETRA.ScanDescriptionFactory.fromZarr} or
any \code{PETRA} module function given the store; \code{zarrToBeamtimeH5.py} restores the HDF5 file.

\end{document}
